\chapter{Contexte : l'apprentissage par renforcement et les seize failles}\label{chap:02-contexte}

Ce chapitre pose les notions dont dépend la suite du rapport : le cadre formel de l'apprentissage par
renforcement (section~\ref{sec:cadre}), les quatre programmes de nos expériences et l'algorithme TRPO
(section~\ref{sec:ctx-programmes}), les tâches simulées (section~\ref{sec:ctx-taches}), ce qu'une
graine aléatoire fixe et comment on évalue un agent (sections~\ref{sec:ctx-graines}
et~\ref{sec:ctx-evaluation}). Il dresse ensuite l'inventaire des seize failles que Joelle Pineau a
montrées en 2018 (section~\ref{sec:failles}), situe notre travail parmi les travaux voisins
(section~\ref{sec:voisins}) et fixe le vocabulaire employé pour les modèles de langue
(section~\ref{sec:vocabulaire}).

\section{Le cadre formel}\label{sec:cadre}

\paragraph{Le processus de décision markovien}
L'apprentissage par renforcement (RL, pour \emph{reinforcement learning}) étudie un programme, l'agent,
qui agit de façon répétée dans un environnement. Le cadre standard est le processus de décision
markovien \parencite{puterman1994mdp,sutton2018rl}, un quintuplet
\begin{equation}
  \mathcal{M} = (\mathcal{S}, \mathcal{A}, P, r, \gamma),
\end{equation}
où $\mathcal{S}$ est l'ensemble des états, $\mathcal{A}$ l'ensemble des actions,
$P(s' \mid s, a)$ la probabilité de passer à l'état $s'$ après l'action $a$ dans l'état $s$,
$r(s, a)$ la récompense reçue et $\gamma \in [0, 1)$ le facteur d'actualisation, qui pèse les
récompenses futures. Une distribution $\rho_0$ tire l'état initial. Dans les tâches de ce rapport,
$\mathcal{S}$ et $\mathcal{A}$ sont des sous-ensembles de $\mathbb{R}^d$ : l'état décrit les positions
et les vitesses des articulations d'un corps simulé, l'action les couples appliqués aux moteurs.

\paragraph{Politique, épisode et retour}
Une politique $\pi_\theta(a \mid s)$ donne la loi de l'action dans chaque état. En RL profond, elle
est un réseau de neurones de paramètres $\theta$. Un épisode est une suite
$\tau = (s_0, a_0, r_0, s_1, a_1, r_1, \dots, s_T)$ obtenue en tirant $s_0 \sim \rho_0$, puis
$a_t \sim \pi_\theta(\cdot \mid s_t)$ et $s_{t+1} \sim P(\cdot \mid s_t, a_t)$, jusqu'à un état
terminal ou jusqu'à une durée maximale. Le retour actualisé d'un épisode et l'objectif de
l'apprentissage s'écrivent
\begin{equation}
  G(\tau) = \sum_{t=0}^{T-1} \gamma^t \, r(s_t, a_t),
  \qquad
  J(\theta) = \mathbb{E}_{\tau \sim \pi_\theta}\bigl[ G(\tau) \bigr].
\end{equation}
L'apprentissage cherche $\theta$ qui maximise $J$. Le score rapporté dans les articles et dans nos
expériences est presque toujours le retour non actualisé $\sum_t r(s_t, a_t)$ d'un épisode, en moyenne
sur plusieurs épisodes d'évaluation.

\paragraph{Fonctions de valeur}
La valeur d'un état et la valeur d'une action sous la politique $\pi$ sont
\begin{equation}
  V^\pi(s) = \mathbb{E}_\pi\Bigl[ \sum_{k \geq 0} \gamma^k r_{t+k} \Bigm| s_t = s \Bigr],
  \qquad
  Q^\pi(s, a) = \mathbb{E}_\pi\Bigl[ \sum_{k \geq 0} \gamma^k r_{t+k} \Bigm| s_t = s, a_t = a \Bigr].
\end{equation}
L'avantage $A^\pi(s, a) = Q^\pi(s, a) - V^\pi(s)$ mesure ce qu'une action rapporte de plus que la
moyenne des actions de la politique dans cet état. Un algorithme acteur-critique apprend à la fois
une politique (l'acteur) et une estimation d'une fonction de valeur (le critique), qui sert à réduire
la variance du signal d'apprentissage de l'acteur.

\paragraph{Gradient de politique}
Pour une politique dérivable en $\theta$, le théorème du gradient de politique \parencite{sutton2018rl}
donne
\begin{equation}
  \nabla_\theta J(\theta) = \mathbb{E}_{\pi_\theta}\bigl[ \nabla_\theta \log \pi_\theta(a_t \mid s_t)\, A^{\pi_\theta}(s_t, a_t) \bigr].
\end{equation}
Les méthodes dites \emph{on-policy}, comme TRPO et PPO, estiment cette espérance avec des épisodes
tirés de la politique courante, puis jettent ces épisodes après la mise à jour. Les méthodes
\emph{off-policy}, comme SAC et CrossQ, gardent les transitions passées dans une mémoire de rejeu et
apprennent un critique $Q$ à partir de ces transitions. Elles réutilisent chaque transition plusieurs
fois. Pour un même nombre de pas d'environnement, elles apprennent en général plus vite, au prix d'un
calcul plus lourd par pas.

\paragraph{Hyperparamètres}
Les paramètres $\theta$ s'apprennent. D'autres réglages restent fixés par l'expérimentateur : le taux
d'apprentissage, $\gamma$, la taille des réseaux, la taille des lots, le nombre de pas collectés avant
chaque mise à jour. Ce sont les hyperparamètres. Les failles F4 à F7 de la section~\ref{sec:failles} les
concernent.

\section{Les quatre programmes et TRPO}\label{sec:ctx-programmes}

Nos expériences emploient quatre programmes, c'est-à-dire quatre implémentations précises d'un
algorithme, dans une bibliothèque et une version données. Le mot \enquote{algorithme} désigne la
méthode décrite dans un article. Le mot \enquote{programme} désigne le code qui l'exécute. Deux
programmes du même algorithme donnent parfois des résultats différents : c'est la faille F14. Les
versions figées de chaque bibliothèque sont données au tableau~\ref{tab:env}. Le chapitre~\ref{chap:05-experiences-methode}
détaille les budgets et l'évaluation.

\paragraph{PPO, en PyTorch et en JAX}
PPO (\emph{Proximal Policy Optimization}) \parencite{schulman2017ppo} est un algorithme acteur-critique
\emph{on-policy}. Avec le rapport de probabilités
$\rho_t(\theta) = \pi_\theta(a_t \mid s_t) / \pi_{\theta_\text{old}}(a_t \mid s_t)$ et un estimateur
$\hat{A}_t$ de l'avantage, il maximise l'objectif tronqué
\begin{equation}
  L^{\text{CLIP}}(\theta) = \hat{\mathbb{E}}_t\Bigl[ \min\bigl( \rho_t(\theta)\, \hat{A}_t,\;
  \operatorname{clip}\bigl(\rho_t(\theta), 1 - \epsilon, 1 + \epsilon\bigr)\, \hat{A}_t \bigr) \Bigr],
\end{equation}
auquel s'ajoutent une perte sur le critique et un bonus d'entropie. La troncature supprime l'intérêt, pour une mise à
jour, d'éloigner trop la nouvelle politique de l'ancienne. Le programme \algo{PPO PyTorch} est celui de
Stable-Baselines3 \parencite{raffin2021sb3}. Le programme \algo{PPO JAX} est celui de sbx, la réécriture
de Stable-Baselines en JAX \parencite{raffin2023sbx,bradbury2018jax}. Les deux programmes reçoivent
les mêmes hyperparamètres d'algorithme et collectent l'expérience par paquets de \num{2048} pas. Leur
initialisation par défaut diffère. Stable-Baselines3 active l'option \file{ortho\_init} : les couches
cachées reçoivent une initialisation orthogonale de gain $\sqrt{2}$, la sortie de l'action un gain de
0,01 et celle de la valeur un gain de 1 (\file{stable\_baselines3/common/policies.py}). La bibliothèque sbx la
désactive par défaut et garde l'initialisation de Flax. Activée, l'option de sbx ne touche que la couche de
sortie de l'action (\file{sbx/ppo/policies.py}). L'expérience E3 compare ces deux programmes
(chapitre~\ref{chap:06-experiences-resultats}). PPO était la référence de 2018. Le protocole le retient parce qu'il reste une
référence courante en 2026 et que GRPO, l'algorithme de notre volet sur les modèles de langue
(chapitre~\ref{chap:07-modeles-de-langue}), en dérive.

\paragraph{SAC}
SAC (\emph{Soft Actor-Critic}) \parencite{haarnoja2018sac} est un algorithme acteur-critique
\emph{off-policy}. Il maximise le retour augmenté de l'entropie de la politique,
\begin{equation}
  J_{\text{SAC}}(\pi) = \sum_t \mathbb{E}_{(s_t, a_t) \sim \pi}\bigl[ r(s_t, a_t) + \alpha\, \mathcal{H}\bigl(\pi(\cdot \mid s_t)\bigr) \bigr],
\end{equation}
où le coefficient $\alpha$ règle l'exploration et s'ajuste pendant l'entraînement. Deux critiques $Q$
appris en parallèle limitent la surestimation des valeurs. Des copies retardées de ces critiques, les
réseaux cibles, stabilisent les cibles d'apprentissage. Le programme \algo{SAC} de nos expériences est
celui de sbx. SAC sert de référence \emph{off-policy} courante en contrôle continu.

\paragraph{CrossQ}
CrossQ \parencite{bhatt2024crossq} descend de SAC. Il retire les réseaux cibles et normalise les
activations par lots. Les paires $(s, a)$ et $(s', a')$ passent ensemble dans le critique, pour que la
normalisation voie les deux distributions. Ses auteurs annoncent une efficacité en données au moins
égale à celle de méthodes beaucoup plus coûteuses en calcul. Le programme \algo{CrossQ} de nos
expériences est celui de sbx. Il emploie une renormalisation par lots (\emph{batch renormalization}, \cite{ioffe2017batchrenorm}),
dans l'acteur comme dans le critique, avec un élan de 0,99. Pendant les \num{100000} premières mises à jour
de chaque réseau, la renormalisation se comporte comme une normalisation par lots ordinaire. Rapporté à
notre budget de pas, ce seuil laisse le critique dans ce régime pendant la première moitié de
l'entraînement. L'acteur, mis à jour une fois sur trois, n'atteint jamais le seuil
(section~\ref{sec:programmes}). Son optimiseur
Adam prend $\beta_1 = 0{,}5$ au lieu de 0,9 (\file{sbx/crossq/policies.py}). Le tableau~\ref{tab:hyperparametres}
donne ses autres réglages.

\paragraph{TRPO}
TRPO (\emph{Trust Region Policy Optimization}) \parencite{schulman2015trpo} est l'algorithme de la
figure la plus célèbre de 2018 (faille F11). Il maximise un objectif de substitution sous une
contrainte de divergence de Kullback-Leibler,
\begin{equation}
  \max_\theta\; \hat{\mathbb{E}}_t\bigl[ \rho_t(\theta)\, \hat{A}_t \bigr]
  \quad \text{sous} \quad
  \hat{\mathbb{E}}_t\bigl[ \mathrm{KL}\bigl( \pi_{\theta_\text{old}}(\cdot \mid s_t) \,\big\|\, \pi_\theta(\cdot \mid s_t) \bigr) \bigr] \leq \delta,
\end{equation}
résolue par un gradient conjugué suivi d'une recherche linéaire. PPO en est une simplification. Nous
l'employons seulement dans l'expérience B1, avec le programme de sb3-contrib~2.9.2 et ses valeurs par
défaut, sans aucun réglage (chapitre~\ref{chap:05-experiences-methode}).

\section{Les tâches MuJoCo}\label{sec:ctx-taches}

Henderson et ses coauteurs comparaient leurs algorithmes sur des tâches du simulateur MuJoCo
(\emph{Multi-Joint dynamics with Contact}) \parencite{todorov2012mujoco,henderson2018matters}. Nous
prenons les mêmes tâches dans leur version de 2026, celle de Gymnasium \parencite{towers2024gymnasium}.
Le tableau~\ref{tab:taches} donne les dimensions des espaces d'état et d'action, lues dans
l'environnement figé de nos expériences. Chaque épisode dure au plus \num{1000} pas.
\task{Hopper} et \task{Walker2d} s'arrêtent plus tôt quand le corps tombe.

\begin{table}[tbp]
\centering\small
\caption{Les quatre tâches de nos expériences, dans Gymnasium~1.1.1 et MuJoCo. Les dimensions et la
durée maximale d'un épisode sont lues dans l'environnement de nos expériences (\file{sim/env}).}
\label{tab:taches}
\begin{tabular}{@{}l r r >{\raggedright\arraybackslash}p{7.4cm}@{}}
\toprule
Tâche & état & action & Corps simulé et récompense \\
\midrule
\task{HalfCheetah-v5} & 17 & 6 & guépard articulé en deux dimensions. Récompense : vitesse vers l'avant, moins un coût des commandes \\
\task{Hopper-v5} & 11 & 3 & sauteur à une jambe. Récompense : vitesse vers l'avant, bonus tant que le corps reste debout \\
\task{Walker2d-v5} & 17 & 6 & marcheur à deux jambes. Récompense : mêmes termes que \task{Hopper} \\
\task{Swimmer-v5} & 8 & 2 & nageur à trois segments dans un fluide visqueux. Récompense : vitesse vers l'avant \\
\bottomrule
\end{tabular}
\end{table}

Pineau jugeait ces bancs d'essai \enquote{ridiculement simples} : des états décrits par peu de
nombres, des transitions déterministes, un monde qui pèse moins de 100~kilo-octets
\parencite{pineau2018neurips}. Cette critique fonde la faille F16. Nous gardons ces tâches, puisque le
but est de refaire les expériences de 2018 telles qu'elles furent faites.

\section{Ce que change une graine aléatoire}\label{sec:ctx-graines}

\paragraph{Définition}
Un entraînement tire beaucoup de nombres au hasard. Un ordinateur fabrique ce hasard à partir d'un
entier de départ, la graine aléatoire (\emph{random seed}). Dans nos programmes, la graine aléatoire est
passée à la bibliothèque (\file{seed=SEED} dans \file{sim/run.py}). La bibliothèque en dérive l'initialisation
des poids des réseaux, le tirage des actions d'exploration, l'ordre des mini-lots et la suite des états
initiaux de l'environnement d'entraînement. Dans ce rapport, un \emph{essai} désigne un entraînement
complet avec une graine aléatoire donnée.

\paragraph{Ce qui reste fixe}
Même programme, même graine aléatoire, mêmes versions des bibliothèques et même machine : le protocole
attend le même résultat. Nos journaux le vérifient pour \algo{PPO PyTorch} et \algo{PPO JAX} : un essai
arrêté puis relancé retrouve à l'identique ses points d'évaluation. Pour \algo{SAC} et \algo{CrossQ},
nous ne l'avons pas vérifiée (chapitre~\ref{chap:10-limites}).

\paragraph{Ce qui change}
Changer la graine aléatoire change l'initialisation des poids, l'exploration, l'ordre des mini-lots et
les états initiaux. Le score final d'un même programme varie alors
beaucoup. Sur nos trente essais de \algo{PPO PyTorch} sur \task{HalfCheetah}, l'écart type vaut 34~\% de
la moyenne (chapitre~\ref{chap:06-experiences-resultats}).

\paragraph{L'unité statistique est l'essai}
La graine aléatoire est une source de variance au même titre que le tirage d'un échantillon dans une
enquête. Une comparaison entre deux programmes repose donc sur plusieurs essais indépendants par
programme, avec une incertitude calculée sur ces essais, sinon elle mesure le hasard du tirage.
L'épisode ou le point d'une courbe ne compte pas pour une observation indépendante, puisqu'il vient
d'un seul essai. La faille F11 et l'expérience B1 montrent ce qui arrive quand on prend l'épisode pour
l'unité (chapitre~\ref{chap:06-experiences-resultats}). Le critère X8 du sondage demande, pour cette
raison, si les articles calculent leur incertitude sur des essais indépendants
(section~\ref{sec:sm-x8}).

\section{Évaluer un agent}\label{sec:ctx-evaluation}

\paragraph{Le score d'un essai}
Évaluer un agent revient à jouer des épisodes avec sa politique, sans apprentissage, puis à faire la
moyenne de leurs retours. Trois choix restent à faire. Un article qui ne les dit pas montre la faille F8.
\begin{itemize}
  \item Le point de mesure : le score à la fin de l'entraînement, le meilleur score atteint le long de
  la courbe, ou la moyenne de la courbe, qui approche l'aire sous la courbe divisée par sa longueur.
  \item Le mode de la politique : stochastique comme à l'entraînement, ou déterministe (l'action la
  plus probable).
  \item L'environnement : celui de l'entraînement, ou des conditions tenues à part, comme d'autres
  graines aléatoires d'environnement ou une physique modifiée (faille F15).
\end{itemize}
Nos programmes sont évalués tous les \num{50000} pas sur 10 épisodes, en mode déterministe, avec des
graines aléatoires d'environnement $\num{10000} + i$, où $i$ numérote les épisodes d'évaluation, qui ne servent
jamais à l'entraînement (chapitre~\ref{chap:05-experiences-methode}).

\paragraph{Agréger les essais}
Avec $n$ essais de scores $x_{(1)} \leq \dots \leq x_{(n)}$, la moyenne et la médiane réagissent
différemment aux essais extrêmes. \textcite{agarwal2021precipice} recommandent la moyenne interquartile
(IQM, \emph{interquartile mean}), qui écarte le quart le plus bas et le quart le plus haut des essais
puis fait la moyenne du reste :
\begin{equation}
  \mathrm{IQM}(x) = \frac{1}{n - 2k} \sum_{i = k + 1}^{n - k} x_{(i)}, \qquad k = \lfloor n/4 \rfloor .
\end{equation}
C'est la moyenne tronquée à 25~\% de \file{scipy.stats.trim\_mean}, que la bibliothèque rliable
appelle : avec dix essais, l'IQM est la moyenne des six essais du milieu. Son intervalle de confiance
se calcule par un bootstrap stratifié : on rééchantillonne les essais avec remise au sein de chaque
tâche, on recalcule l'IQM, puis on prend les quantiles à 2,5~\% et 97,5~\% des \num{2000} répliques
(\file{src/pineau\_rl/bootstrap.py}).

\paragraph{Comparer deux programmes}
La probabilité d'amélioration de $X$ sur $Y$ est la probabilité qu'un essai de $X$ batte un essai de
$Y$ tiré indépendamment. Avec $n_X$ essais de $X$ et $n_Y$ essais de $Y$, on l'estime par
\begin{equation}
  P(X > Y) = \frac{1}{n_X n_Y} \sum_{i=1}^{n_X} \sum_{j=1}^{n_Y} \Bigl( \mathbb{1}[x_i > y_j] + \tfrac{1}{2}\, \mathbb{1}[x_i = y_j] \Bigr).
\end{equation}
Cette estimation vaut la statistique $U$ de Mann-Whitney divisée par $n_X n_Y$. À 0,5, un essai de $X$
bat un essai de $Y$ une fois sur deux. Nos expériences déclarent un écart établi quand l'intervalle
bootstrap à 95~\% de $P(X > Y)$ exclut 0,5, ou quand les intervalles des deux IQM ne se chevauchent
pas : c'est la règle du vainqueur établi. Sur des groupes d'essais d'un même programme, où aucun écart
n'existe, elle conclut à tort un peu plus
d'une fois sur vingt avec dix essais par groupe, près d'une fois sur dix avec cinq
(section~\ref{sec:regle}). Le test $t$ de Welch \parencite{welch1947generalization}, qui compare deux
moyennes sans supposer les variances égales, sert seulement dans les expériences E1 et B1, puisque
c'est lui que la figure de 2018 employait.

\section{Les seize failles montrées en 2018}\label{sec:failles}

Dans ses deux conférences de 2018, à ICLR puis à NeurIPS \parencite{pineau2018iclr,pineau2018neurips},
Pineau montre chaque faille par un exemple, souvent par une figure tirée de l'article de son équipe
signé Peter Henderson et ses coauteurs \parencite{henderson2018matters}. Nous avons relevé seize
failles, dans l'ordre des deux conférences, puis nous les avons rangées en quatre familles. Le tableau~\ref{tab:failles} les liste avec la diapositive ou la page qui les montre et l'instrument de ce
rapport qui les mesure. Le numéro F1 à F16 sert seulement à retrouver une faille dans les tableaux.
L'annexe~\ref{chap:A-grille} donne les règles de décision de chaque critère de lecture.

\begin{small}
\begin{longtable}{@{}>{\raggedright\arraybackslash}p{0.6cm} >{\raggedright\arraybackslash}p{5.3cm} >{\raggedright\arraybackslash}p{3.6cm} >{\raggedright\arraybackslash}p{4.4cm}@{}}
\caption{Les seize failles que Joelle Pineau a montrées en 2018 dans l'apprentissage par renforcement
profond. Pour chacune : la diapositive de sa conférence d'ICLR ou la page de son document de NeurIPS
qui la montre, puis les critères de lecture (Q, X, Y, chapitre~\ref{chap:03-sondage-methode}) et les
expériences (E, B, chapitre~\ref{chap:05-experiences-methode}) qui la mesurent dans ce rapport.}\label{tab:failles}\\
\toprule
& Faille & Où Pineau la montre & Ce qui la mesure ici \\
\midrule
\endfirsthead
\toprule
& Faille & Où Pineau la montre & Ce qui la mesure ici \\
\midrule
\endhead
\bottomrule
\endfoot
\multicolumn{4}{@{}l}{\emph{Le matériel : ce qu'il faut pour refaire}} \\
F1 & L'article publié sans son code & diapositive 6 (ICLR), page 29 (NeurIPS) & critère Q5 \\
F2 & Aucun plan pour refaire l'expérience, comme pour AlphaGo & diapositives 9 à 12 (ICLR) & critères Q3 et X7 \\
F3 & Le matériel et les logiciels, qui font partie du résultat & pages 34 et 35 (NeurIPS) & critères X7 et Y10 \\
\midrule
\multicolumn{4}{@{}l}{\emph{Le réglage : hyperparamètres et références}} \\
F4 & Les hyperparamètres de la méthode, donnés en partie & page 29 (NeurIPS) & critère Q3 \\
F5 & La façon de choisir les hyperparamètres, tue & page 29 (NeurIPS) & critère Q6 \\
F6 & Des réglages qui changent tout et qui interagissent & diapositives 24 et 25 (ICLR) & critère Y9 \\
F7 & Des références moins bien réglées que la méthode & diapositive 25 (ICLR), page 29 (NeurIPS) & critères Q4 et X4, expérience E6 \\
\midrule
\multicolumn{4}{@{}l}{\emph{Les graines aléatoires et la statistique}} \\
F8 & Une mesure de performance jamais définie & diapositives 26 et 27 (ICLR) & critères Y1 et Y2, expérience B3 \\
F9 & Trop peu d'essais : trois ou cinq, jamais dix dans son tableau & diapositive 28 (ICLR) & critère X1, expérience E5 \\
F10 & Les meilleurs essais sur $N$, seuls publiés & diapositives 29 et 30 (ICLR) & critère X3, expériences E1 et E4 \\
F11 & Même code, autre graine aléatoire : deux \enquote{algorithmes} & diapositive 31 (ICLR), page 22 (NeurIPS) & expériences E1 et B1 \\
F12 & Aucun test statistique & page 29 (NeurIPS) & critères Q8, X2, Y11 et Y12 \\
\midrule
\multicolumn{4}{@{}l}{\emph{La portée des conclusions}} \\
F13 & Le meilleur algorithme qui change avec la tâche & diapositives 20 et 21 (ICLR) & critères Y3 et Y4, expérience E2 \\
F14 & Deux implémentations du même algorithme qui divergent & diapositives 22 et 23 (ICLR) & critère Y5, expérience E3 \\
F15 & L'évaluation sur les conditions d'entraînement & pages 36 à 51 (NeurIPS) & critères Q7 et Y8, expérience E7 \\
F16 & Des bancs d'essai trop simples, faciles à apprendre par cœur & pages 43 à 50 (NeurIPS) & critères X6 et Y7 \\
\end{longtable}
\end{small}

\subsection{Le matériel : ce qu'il faut pour refaire}

\paragraph{F1, l'article publié sans son code}
Un article de RL décrit sa méthode en quelques pages. Le programme qui a produit les résultats compte
des milliers de lignes, pleines de choix que la prose ne mentionne pas. Pineau reprend la formule de
\textcite{buckheit1995wavelab} citée au chapitre~\ref{chap:01-introduction}, où ils condensaient une idée de Jon Claerbout :
l'article est la publicité, la recherche est le code et les instructions qui ont produit les figures.
Sur les 50 articles de RL de 2018 qu'elle a lus, 55~\% donnent un lien vers leur code.

\paragraph{F2, aucun plan pour refaire l'expérience}
Le code ne suffit pas sans les réglages exacts et une idée du calcul nécessaire. L'exemple de Pineau
est AlphaGo \parencite{silver2016alphago}. L'article est détaillé, avec une longue annexe. Sa
diapositive demande pourtant si cela suffit : il faudrait aussi les parties d'experts, le code, les
hyperparamètres, les ordinateurs et le protocole d'évaluation. Une découverte qu'aucun laboratoire
extérieur n'a les moyens de refaire reste une affirmation.

\paragraph{F3, le matériel et les logiciels}
Le même code lancé sur deux machines donne parfois deux résultats : les cartes graphiques n'arrondissent
pas les calculs comme les processeurs, une nouvelle version d'une bibliothèque change une valeur par
défaut. La liste de contrôle de Pineau demande l'infrastructure de calcul employée et le code
\enquote{avec toutes ses dépendances}.

\subsection{Le réglage : hyperparamètres et références}

\paragraph{F4, les hyperparamètres de la méthode, donnés en partie}
Un seul réglage oublié suffit à empêcher la reproduction : le lecteur prend la valeur par défaut de sa
bibliothèque et obtient autre chose. Pineau compte que 90~\% des articles de 2018 donnent tous les
hyperparamètres de leur méthode.

\paragraph{F5, la façon de choisir les hyperparamètres, tue}
Une valeur finale ne dit pas comment on l'a trouvée. Une équipe qui a essayé cent combinaisons et
gardé la meilleure a optimisé sur son banc d'essai. Le lecteur a besoin du budget de recherche. Seuls
20~\% des articles de 2018 le disent.

\paragraph{F6, des réglages qui changent tout}
Henderson et ses coauteurs changent un réglage à la fois. Pour PPO sur \task{HalfCheetah}, passer d'un
réseau de deux couches de 64 neurones à un réseau de 400 puis 300 neurones fait tomber le score de
\num{2201} à $\num{-1180}$. Les réglages interagissent : le meilleur réseau pour un algorithme cesse de l'être
pour un autre. \textcite{engstrom2020implementation} puis \textcite{andrychowicz2021matters} ont confirmé
cette sensibilité à grande échelle.

\paragraph{F7, des références moins bien réglées que la méthode}
L'auteur d'une méthode nouvelle la règle avec soin. Il prend souvent la référence avec ses valeurs par
défaut. La comparaison oppose alors une méthode réglée à une méthode qui ne l'est pas. Pineau compte
que 60~\% des articles de 2018 donnent les hyperparamètres de leurs références.
\textcite{eimer2023hyperparameters} réclament le même budget de réglage pour les deux.

\subsection{Les graines aléatoires et la statistique}

\paragraph{F8, une mesure de performance jamais définie}
Le score final, la moyenne sur l'entraînement et le meilleur score atteint ne classent pas toujours les
méthodes dans le même ordre (section~\ref{sec:ctx-evaluation}). Un article qui ne dit pas laquelle il emploie, ni
ce que représentent ses barres d'erreur, ne se compare à aucun autre.

\paragraph{F9, trop peu d'essais}
Pineau montre un tableau d'articles de 2015 à 2017 : ils lancent trois ou cinq graines aléatoires,
aucun n'atteint dix. Avec trois essais, une moyenne dépend beaucoup du hasard d'un seul essai.
\textcite{colas2018seeds} calculent le nombre de graines aléatoires nécessaire pour détecter un écart donné.

\paragraph{F10, les meilleurs essais sur $N$}
Une équipe lance dix essais, garde les meilleurs et publie leur moyenne. Chaque chiffre est vrai, mais
la moyenne publiée dépasse ce que la méthode donne en moyenne. La dispersion paraît aussi plus petite
qu'elle n'est. Pineau l'illustre par dix courbes dont elle garde les trois meilleures.

\paragraph{F11, même code, autre graine aléatoire}
L'équipe de Henderson entraîne dix fois TRPO sur \task{HalfCheetah} en changeant seulement la
graine aléatoire, forment deux groupes de cinq et tracent les deux courbes moyennes. Elles sont
nettement séparées. Un test les déclare différentes ($t = -9{,}09$, $p = 0{,}0016$). Pineau précise
sur scène que les groupes ne sont pas triés. Les chapitres~\ref{chap:05-experiences-methode} et~\ref{chap:06-experiences-resultats}
reviennent longuement sur ce chiffre (expériences E1 et B1).

\paragraph{F12, aucun test statistique}
Quand deux méthodes diffèrent de quelques points, un test ou un intervalle dit si l'écart dépasse ce
que le hasard des graines aléatoires produit seul. Dans la lecture de Pineau, 5~\% des articles de
2018 appliquent un test à leurs comparaisons.

\subsection{La portée des conclusions}

\paragraph{F13, le meilleur algorithme qui change avec la tâche}
L'article de Henderson classe quatre algorithmes sur quatre tâches. L'algorithme en tête sur
\task{HalfCheetah} cesse de l'être sur \task{Hopper} ou \task{Swimmer}. Une conclusion générale tirée
d'une ou deux tâches généralise un cas particulier.

\paragraph{F14, deux implémentations qui divergent}
Trois implémentations de TRPO, avec le même réseau, finissent sur \task{HalfCheetah} à \num{1939}, 834 et
205. Comparer sa méthode à \enquote{TRPO} sans dire quelle implémentation on a prise revient à la
comparer à un programme anonyme.

\paragraph{F15, l'évaluation sur les conditions d'entraînement}
En apprentissage supervisé, on évalue sur des données jamais vues. En RL, l'agent s'évalue souvent dans
le monde même de son entraînement. Pineau projette une phrase qu'elle met en débat, selon laquelle le
RL serait le seul cas de l'apprentissage automatique où tester sur son ensemble d'entraînement est
acceptable. Elle propose de séparer les graines aléatoires d'entraînement et de test, puis les tâches
\parencite{zhang2018dissection,zhang2018natural}. Pineau ne trouve que 10~\% d'articles de 2018 qui
évaluent sur des conditions tenues à part.

\paragraph{F16, des bancs d'essai trop simples}
Un agent assez gros retient la bonne suite d'actions d'un banc d'essai déterministe sans rien
comprendre du monde \parencite{machado2018revisiting}. Pineau réclame aussi la version de chaque banc
d'essai, pour qu'un résultat se refasse dans le même monde.

\paragraph{Les remèdes de 2018}
Aux seize failles, Pineau oppose d'abord une liste de contrôle de reproductibilité (pages 31 à
33 à NeurIPS). Pour chaque algorithme : une description claire, une analyse de complexité, le lien vers
le code et ses dépendances. Pour chaque résultat théorique : l'énoncé, les hypothèses, la preuve. Pour
chaque figure : la collecte des données, la répartition entre entraînement, validation et test, les
données exclues, la plage et la méthode de choix des hyperparamètres, le nombre exact d'essais, la
définition des statistiques et des barres d'erreur, le matériel employé. Pineau y ajoute le défi de
reproduction d'ICLR 2018 (diapositives 37 à 43 à ICLR), présenté au chapitre~\ref{chap:01-introduction}.
NeurIPS a adopté la liste en 2019 \parencite{pineau2021improving}.

\section{Travaux voisins}\label{sec:voisins}

Au début du projet, nous avons cherché sur arXiv, sur OpenReview et dans les actes des congrès. À notre
connaissance, aucun travail ne réapplique la grille de Pineau à des articles de RL récents en la
calibrant sur 2018. Aucun ne refait le protocole de Henderson avec les algorithmes de
2026 sous un protocole écrit d'avance. Les travaux voisins se rangent en cinq groupes.

\paragraph{L'expérience de 2018}
\textcite{henderson2018matters} fondent la plupart des failles F6 à F14 : sensibilité aux
hyperparamètres, aux implémentations, aux graines aléatoires, aux tâches. \textcite{islam2017reproducibility}
les précèdent sur les mêmes tâches. \textcite{lynnerup2020survey} refont des algorithmes sur un vrai
bras robotique. \textcite{bjorck2022variance} montrent, pour le contrôle à partir d'images, qu'une part
de la variance entre graines aléatoires vient d'essais qui échouent tôt par instabilité numérique.

\paragraph{La méthode statistique}
\textcite{colas2018seeds} font une analyse de puissance du nombre de graines aléatoires. \textcite{jordan2020evaluating}
discutent la mesure de performance elle-même, puis \textcite{jordan2024position} le coût d'une
comparaison rigoureuse. \textcite{agarwal2021precipice} apportent l'IQM, les intervalles par bootstrap
stratifié, les profils de performance et la probabilité d'amélioration, réunis dans la bibliothèque
rliable. Ils montrent que plusieurs améliorations annoncées sur les jeux Atari se fondent dans
l'incertitude une fois les intervalles calculés. \textcite{patterson2024empirical} en tirent un guide.
\textcite{eimer2023hyperparameters} réclament un budget de réglage égal pour les références. Ces travaux
disent quoi faire. Notre sondage mesure si les articles le font.

\paragraph{Les listes de contrôle et les défis de reproduction}
\textcite{pineau2021improving} décrivent le programme de reproductibilité de NeurIPS 2019 : la liste de
contrôle obligatoire, le défi de reproduction et la politique de dépôt du code. La part d'articles
acceptés avec du code passe de moins de 50~\% à 74,4~\%. \textcite{magnusson2023nlpchecklist} mesurent
l'effet de la liste de contrôle des congrès de traitement automatique des langues.

\paragraph{Les audits de la littérature}
\textcite{gundersen2018state} notent 400 articles des congrès IJCAI et AAAI sur la documentation des
méthodes, des données et des expériences. \textcite{coakley2026shift} suivent \num{56800} articles de cinq
congrès d'intelligence artificielle de 2014 à 2024 : la part qui partage code et données passe de 11~\% à 64~\%. La hausse commence avant les listes de contrôle. En RL, l'audit le plus proche est celui de
\textcite{gorsane2022marl} : 75 articles de RL coopératif à plusieurs agents, de 2016 à 2022, dont
environ un tiers ne quantifie aucune incertitude. Ces audits laissent de côté le RL à un seul agent et
la grille de 2018.

\paragraph{Les modèles de langue}
\textcite{hochlehnert2025sober} réévaluent des modèles de raisonnement et trouvent des écarts types de 5
à 15 points de pourcentage d'une graine aléatoire d'échantillonnage à l'autre sur de petits bancs
d'essai comme AIME~2024. \textcite{chandak2025incorrect} relèvent, dans un billet non relu par des pairs,
des modèles de départ sous-évalués dans sept articles, l'équivalent moderne de la faille F7.
\textcite{shao2025spurious} montrent qu'une récompense tirée au sort améliore les modèles Qwen en
mathématiques, mais pas les modèles Llama ou OLMo. \textcite{yue2025rlvr} montrent que la mesure pass@$k$
renverse un classement selon $k$, un exemple de la faille F8.

\section{Le vocabulaire des modèles de langue}\label{sec:vocabulaire}

Le vocabulaire du RL a gagné les modèles d'IA générative qui écrivent du texte. On emprunte au RL
certains de ses algorithmes, avec une ressemblance vague, pour affiner les réponses d'un modèle de
langue préentraîné, sans agent qui agit dans un monde. Nous fixons nos mots comme suit.

\begin{itemize}
  \item \textbf{RL classique.} Le RL au sens de la section~\ref{sec:cadre} : un agent, souvent entraîné de
  zéro, agit dans un environnement doté d'états qui évoluent, en robotique, dans les jeux ou en
  contrôle. Le RL hors ligne en fait partie. C'est la famille \file{classic} du critère X5.
  \item \textbf{RL de grands modèles.} Une étiquette de commodité pour les articles qui appellent RL ce
  qu'ils appliquent à un modèle génératif préentraîné conditionné par le langage (modèle de langue,
  modèle de vision-langage, modèle texte-image). C'est la famille \file{llm} du critère X5
  (chapitre~\ref{chap:03-sondage-methode}). L'étiquette reprend le mot des articles. Elle ne vaut pas jugement
  sur la nature de ce qu'ils font.
  \item \textbf{RLHF.} L'apprentissage à partir de retours humains (\emph{reinforcement learning from
  human feedback}) \parencite{christiano2017deep,ouyang2022instructgpt}. Des annotateurs classent des
  réponses, un modèle de récompense apprend leurs préférences, puis le modèle de langue est affiné pour
  maximiser ce score. Le sigle contient les lettres RL. Nous ne rangeons pas pour autant RLHF dans
  le RL classique : la politique optimisée produit un texte en réponse à une invite, puis reçoit une seule note,
  sans environnement dont l'état évolue sous ses actions.
  \item \textbf{RLVR et GRPO.} Le RLVR (\emph{reinforcement learning with verifiable rewards}) remplace
  le modèle de récompense par un programme qui vérifie la réponse : la réponse finale d'un problème de
  mathématiques comparée à la bonne, un code soumis à des tests \parencite{deepseek2025r1}. GRPO
  (\emph{Group Relative Policy Optimization}) \parencite{shao2024deepseekmath} est l'algorithme le plus
  répandu de cette famille. Il dérive de PPO : il tire un groupe de réponses pour chaque invite et
  remplace le critique par la récompense de chaque réponse centrée et réduite au sein du groupe. Notre
  volet GRPO est décrit au chapitre~\ref{chap:07-modeles-de-langue}.
\end{itemize}

Les grands modèles de langue sont, dans ce rapport, des modèles d'IA générative. Les lecteurs de notre
sondage sont des agents bâtis sur un tel modèle (chapitre~\ref{chap:03-sondage-methode}). Le mot \enquote{agent}
désigne alors un logiciel qui lit un PDF et remplit une grille, pas l'agent d'un processus de décision
markovien.
